%% 1. Introduction  [~2 pages]
\chapter{Introduction}\label{ch:intro}

% - Contexte : place dans le cursus (M2 Info Montpellier, alternance),
%   entreprise (Namirial SPA), dates.
% - Présentation succincte des missions.
% - Notions techniques nécessaires à la compréhension de la suite.
% - Annonce du plan.

\section*{Contexte}

Ce mémoire rend compte de mon alternance de Master~2 Informatique à l'Université de
Montpellier (Faculté des Sciences), effectuée au sein de \textbf{Namirial SPA}, dans
la division \textbf{Identité \& Onboarding}.
% À COMPLÉTER : dates précises de l'alternance, rythme, équipe de rattachement.

\section*{Présentation des missions}

Au cours de cette alternance, j'ai contribué à deux missions principales, complétées
par plusieurs missions annexes. Les deux missions principales s'inscrivent dans le
traitement automatisé des pièces d'identité. La première relève de la \textbf{lutte
contre la fraude documentaire}, la seconde de l'\textbf{extraction d'informations}.
Les missions annexes, plus ponctuelles, m'ont permis d'aborder d'autres facettes du
traitement intelligent de documents, à savoir l'évaluation de modèles d'OCR de
nouvelle génération, l'optimisation d'un service interne d'extraction et un projet
client de \emph{débundling} de documents.

Toutes partageaient un même contexte, celui des produits de vérification d'identité.
Les modèles d'apprentissage automatique n'y sont pas une fin en soi. Ils doivent être
directement industrialisés dans la plateforme et satisfaire des exigences fortes de
fiabilité, de robustesse et de conformité.

\section*{Notions techniques nécessaires}

La compréhension des chapitres suivants suppose quelques notions de base.

\begin{itemize}
  \item les \textbf{réseaux de neurones convolutifs} (CNN) pour la classification
        d'images, et en particulier les architectures \textbf{ResNet} (connexions
        résiduelles) et \textbf{ShuffleNet}, sa variante allégée ;
  \item les notions d'\textbf{ingénierie des caractéristiques} (\emph{feature
        engineering}) et de représentation de la couleur (espaces colorimétriques) ;
  \item les principes de l'\textbf{OCR} et de l'\textbf{extraction d'informations
        clés} (\emph{Key Information Extraction}) ;
  \item l'usage de \textbf{grands modèles de langage} pour l'analyse de documents et
        la \textbf{contrainte de leur sortie} par un schéma typé ou une énumération ;
  \item les techniques d'\textbf{augmentation de données} (\emph{data augmentation}),
        dont l'\emph{inpainting}, face à un volume de données limité.
\end{itemize}

\section*{Organisation du mémoire}

Après une présentation de l'entreprise, le chapitre \textbf{Environnement technique}
pose les bases sur lesquelles reposent les projets, à savoir les langages, les
bibliothèques et les outils internes mobilisés. Le chapitre \textbf{Travaux
effectués} constitue le cœur du mémoire et traite chaque projet de bout en bout, de
sa problématique à ses résultats. Le chapitre \textbf{Veille technologique} relate le
suivi continu des évolutions techniques mené tout au long de l'alternance. La
conclusion dresse enfin un bilan fonctionnel et personnel.
